%===============================================================================
% ifacconf.tex 2022-02-11 jpuente  
% 2022-11-11 jpuente change length of abstract
% Template for IFAC meeting papers
% Copyright (c) 2022 International Federation of Automatic Control
%===============================================================================
\documentclass{ifacconf}

% \usepackage{hyperref}
\usepackage{graphicx}      % include this line if your document contains figures
\usepackage{natbib}        % required for bibliography
\usepackage{comment}
\usepackage{amsmath}
\usepackage{amssymb}
\usepackage[dvipsnames]{xcolor}

\usepackage{subfig}
\usepackage{kotex}
\usepackage{tikz}

\usepackage{lipsum} % for dummy text
% ============================================
%         GENERAL MACROS
% ============================================
\newcommand{\template}{template}
\newcommand{\KHCH}[1]{{\color{RoyalPurple} [KH: #1]}} % Kyunghwan's
\newcommand{\JYKI}[1]{{\color{red} [JY: #1]}} % Jiyun's
\newcommand{\MSRY}[1]{{\color{Peach} [MS: #1]}} % Myeongseok's 
\newcommand\eg{\textrm{e.g.,\ }}
\newcommand\ie{\textrm{i.e.,\ }}

\newcommand{\colorLine}[2]{[\tikz[baseline=(current bounding box.base)]{\draw[color=#1,
#2,line width=2pt] (0,3pt) -- (1.5em,3pt);}]}
% colors
% https://latexcolor.com % color list
\definecolor{airforceblue}{rgb}{0.36, 0.54, 0.66}
\definecolor{awesome}{rgb}{1.0, 0.13, 0.32}
\newcommand{\ctxt}[2]{{\color{#1}{#2}\color{black}}}
\newcommand{\cred}[1]{{\ctxt{awesome}{#1}}}
\newcommand{\cblue}[1]{{\ctxt{airforceblue}{#1}}}

% ============================================
%         MATH MACROS
% ============================================
% \input{\template/macros/macros_math.tex}
\newcommand\R{\mathbb{R}}
\newcommand\N{\mathbb{N}}

\newcommand*{\mv}[1]{\boldsymbol{#1}} % my vector
\newcommand*{\mm}[1]{\boldsymbol{#1}} % my matrix
\newcommand*{\norm}[1]{\lVert #1 \rVert} % euclidean norm

\DeclareMathOperator{\myvec}{vec}       % vectorization
\DeclareMathOperator{\myproj}{{Proj}}   % projection
\DeclareMathOperator{\mysat}{{sat}}       % saturation
\DeclareMathOperator{\myrow}{row}       % row of a matrix
\DeclareMathOperator{\mycol}{Col}       % column of a matrix
\DeclareMathOperator{\mysym}{sym}       % symmetric part of a matrix
\DeclareMathOperator{\mydiag}{diag}     % diagonal matrix
\DeclareMathOperator{\mysign}{sign}     % sign function
\DeclareMathOperator{\mypoly}{Poly}     % polynomial

% ============================================
%         MOTOR SYMBOLS
% ============================================
\newcommand*{\ia}{i_\mathrm{\alpha}} % a-phase current
\newcommand*{\ib}{i_\mathrm{\beta}} % b-phase current

% ============================================
%         FRACTIONS
% ============================================
\newcommand\der{\mathrm d}
\newcommand*{\ddt}{
   \frac{\der}{\der t}
}
\newcommand*{\ddtt}{
   \tfrac{\der}{\der t}
}
\newcommand*{\ddttddtt}{
   \tfrac{\der^2}{\der t^2}
}
\newcommand*{\ddfrac}[2]{
   \frac{\der {#1}}{\der {#2}}
}
\newcommand*{\ddtfrac}[2]{
   \tfrac{\der {#1}}{\der {#2}}
}
\newcommand*{\ppfrac}[2]{
   \frac{\partial {#1}}{\partial {#2}}
}
\newcommand*{\pptfrac}[2]{
   \tfrac{\partial {#1}}{\partial {#2}}
}

\input{\template/symbols/symbols_NN.tex}

%===============================================================================
\begin{document}
\begin{frontmatter}

\title{
   \cblue{Contraction Theory Based Control} Under Input Saturation–LMI Approach\thanksref{footnoteinfo}
} 
% Title, preferably not more than 10 words.

\thanks[footnoteinfo]{Sponsor and financial support acknowledgment
goes here. Paper titles should be written in uppercase and lowercase
letters, not all uppercase.}

\author[First]{Myeongseok Ryu} 
\author[First]{Kyunghwan Choi} 
\author[Second]{Sesun You}

\address[First]{
   Korea Advanced Institute of Science and Technology (KAIST), 
   Daejeon 34051, South Korea (e-mail: \{dding\_98,kh.choi\}@kaist.ac.kr)}
% \address[Second]{
%    Korea Advanced Institute of Science and Technology (KAIST), 
%    Daejeon 34051, South Korea (e-mail: @kaist.ac.kr)}
\address[Second]{
   Keimyung University, 
   Daegu 42601, South Korea (e-mail: yousesun@kmu.ac.kr)}

\begin{abstract}                % Abstract of 50--100 words
   Besides the traditional Lyapunov stability theory, contraction theory offers a interesting perspective on the stability of nonlinear systems.
   It allows us analyze the system stability by examining the convergence of trajectories.
   The recent development of the contraction theory based control has introduced computational efficient methods for obtaining contraction metrics using linear matrix inequalities (LMIs).
   In this paper, the contraction theory based control strategy is extended to account for input saturation by incorporating saturation constraints into the LMI framework for contraction metric calculation.
   The proposed method calculates the contraction metrics while ensuring the satisfaction of input saturation constraints and uses as a feedback mechanism to adjust the control inputs.
   The proposed method is demonstrated its effectiveness through numerical simulations via \cred{Lorenz system}.
\end{abstract}

\begin{keyword}
   Contraction theory, input saturation, linear matrix inequalities (LMIs)
\end{keyword}

\end{frontmatter}
%===============================================================================

\section*{Notation}

In this paper, we use the following notation:
\begin{itemize}
   \item $:=$ denotes \textit{defined as}.
   \item $(\cdot^\top)$ denotes the transpose of a matrix or a vector.
   \item $\mv{x}:=[x_i]_{i\in\{1,\cdots,n\}}\in\R^n$ denotes the state vector.
   \item $\mm A:=[a_{ij}]_{i,j\in\{1,\cdots,n\}}\in\R^{n\times n}$ denotes a matrix.
   \item $\lambda_{i}(\mm A),\ i\in\{\max,\min\}$ denotes the maximum and minimum singular value of $\mm A$, respectively.
   \item $\mm I_n$ denotes the identity matrix of size $n$ and $\mm 0_{n\times m}$ denotes the zero matrix of size $n\times m$.
   \item $\mysym$ denotes the symmetric part of a matrix, \ie $\mysym(\mm A):=\tfrac{1}{2}(\mm A+\mm A^\top)$ (see \cite{Tsukamoto:2021ac}).
\end{itemize}

\section{Introduction}

% 미래생각
% \begin{itemize}
%    \item 차량으로
%    \item 상태 제약조건 추가
%    \item 입력 포화 고려
%    \item 카메이터 검증 (횡제어)
%    \begin{itemize}
%       \item CV-STEM 
%       \item sampling
%       \item sampling -> metric (LMI) <- constraint
%       \item DNN -> interpolate
%    \end{itemize}
% \end{itemize}

\subsection{Motivation}

\cred{\lipsum[1-6]}

\cblue{
   \begin{itemize}
      \item 포화에 대한 사전 연구 더해서 비교
      \item S.S.Ge의 논문 참고 
   \end{itemize}
}

\begin{itemize}
   \item Contraction theory introduction
   \item Limitation 1 (computational cost)
   \item Recent development (LMI)
   \item Limitation 2 (input saturation)
   \begin{itemize}
      \item use penalty term to prevent excessive control input
   \end{itemize}
   \item saturation handling in LMI
\end{itemize}

To the best of our knowledge, there has been no previous work that considers the input saturation in LMI frameworks for contraction metrics.
Hence, the main objective of this paper is to extend the contraction theory based control strategy to account for input saturation simultaneously in the LMI framework.

\subsection{Contributions}

The main contributions of this paper are 
\begin{itemize}
   \item The contraction metric is obtained by using linear matrix inequalities (LMIs) which simplifies the calculation of the contraction metric;
   \item Saturation constraint is incorporated into the LMI framework for contraction metric calculation;
   \item ...; and
   \item The effectiveness and the feasibility of the proposed method is validated through numerical simulations via \cred{Lorenz system}.
\end{itemize}


\section{Preliminaries}

Before presenting the proposed method, we first review the contraction theory and related theorems.

\subsection{Contraction Theory} \label{sec:sub:contraction}

Consider that we have the following deterministic nonlinear system:
\begin{equation} \label{eq:pre:sys}
   \ddtt\mv{x}
   =
   \mv{f}(\mv{x},t)
   .
\end{equation}
The contraction theory commences with the definition of \textit{differential dynamics} of the system \eqref{eq:pre:sys} using \textit{differential displacement} $\delta\mv{x}$ \cite[Chap 4]{Kirk:2004aa}, which is defined as follows:
\begin{equation}
   \ddtt{\delta\mv{x}}
   =
   \pptfrac{\mv{f}}{\mv{x}} \delta\mv{x}
   .
\end{equation}
The differential displacement $\delta\mv{x}$ is the infinitesimal displacement at \textit{fixed time}.

The contraction property of the system \eqref{eq:pre:sys} can be analyzed using the Lyapunov function $V:=V(\mv{x},t)=\tfrac{1}{2}\delta\mv{x}^\top \mm{M} \delta\mv{x}$, where $\mm{M}:=\mm{M}(\mv{x},\mv{x}_d,t)\in\R^{n\times n}$ the \textit{contraction metric}.
The contraction metric $\mm{M}$ is symmetric positive definite matrix and satisfies $\overline{m}\mm{I}_n\succeq\mm{M}\succeq\underline{m}\mm{I}_n$ with $\overline{m},\underline{m}\in\R_{>0}$.
For the system \eqref{eq:pre:sys} to be contracting, the following condition must hold:
\begin{equation} \label{eq:pre:contraction:condition}
   \ddtt\mm{M} 
   + 
   2\mysym(
      \mm{M} \pptfrac{\mv{f}}{\mv{x}}
   ) 
   \leq 
   -2\alpha\mm{M}
   ,
\end{equation}
which leads to $\ddtt V \leq -\alpha V$ for some $\alpha\in\R_{>0}$.
In other words, all solutions of the system \eqref{eq:pre:sys} are contracting to a unique solution with an exponential rate $\alpha$.

It is notable that the initial conditions are exponentially forgotten in the contracting system \cite{LOHMILLER:1998aa}.
Therefore, the contraction property is independent of the initial conditions.

\cred{global?? 하다는것도 쓸까?}

\subsection{Partial Contraction Theory} \label{sec:sub:contraction:partial}

On the other hand, the partial contraction theory investigates the contraction of the system using a virtual system.
Following theorem details the partial contraction theory.

\begin{thm} [see \cite{Wang:2004aa}] \label{thm:partial:contraction}
   Consider a nonlinear system of the form 
   $
      \ddtt\mv{x} = \mv{f}(\mv{x},\mv{x},t)
   $
   and the virtual system 
   $
      \ddtt\mv{q} = \mv{f}(\mv{q},\mv{x},t)
   $
   are contracting with respect to $\mv{q}$.
   If a particular solution of virtual system verifies a smooth specific property, then all solutions of the original system verifies this property exponentially.
\end{thm}

\begin{exmp}
   \cred{
   Assume that the possible solutions of the virtual system defined in Theorem \ref{thm:partial:contraction} are $\mv{\xi}_1(t)$ and $\mv{\xi}_2(t)$ which are the solutions of the nonlinear system $\mv{x}$.
   
   }

\end{exmp}

\subsection{Perturbed Contracting System} \label{sec:sub:contraction:perturbed}

The following theorem presents the contraction property of the perturbed system by adding a bounded perturbation $\mv{d}(\mv{x},t)$ to the system \eqref{eq:pre:sys}:
\begin{equation} \label{eq:pre:sys:perturbed}
   \ddtt\mv{x}
   =
   \mv{f}(\mv{x},t) + \mv{d}(\mv{x},t)
   .
\end{equation}

% Steady state error of perturbed contraction system
\begin{thm} [Eq.~15 \cite{LOHMILLER:1998aa}] \label{thm:contraction:perturbed}
   If the system \eqref{eq:pre:sys} is contracting with the contraction metric $\mm{M}$ which satisfies the condition \eqref{eq:pre:contraction:condition}, the path integral $V_l(\mv{q},\delta\mv{q},t)=\int_{\mv{\xi}_1}^{\mv{\xi}_2}\norm{\mm{\Theta}(\mv{q},t)\delta\mv{q}}$, where $\mv{\xi}_1$ and $\mv{\xi}_2$ are the solutions of \eqref{eq:pre:sys} and \eqref{eq:pre:sys:perturbed}, respectively, and $\mv{q}$ is the virtual 
   \MSRY{ADD!!}
   \begin{equation}
      \norm{\mv\xi_1(t)-\mv\xi_2(t)}
      \le
      \tfrac{V_{l}(0)}{\sqrt{\underline{m}}}
      \exp(-\alpha t)
      +
      \tfrac{\overline d}{\alpha}
      \sqrt{\tfrac{\overline m}{\underline{m}}}
      (1-\exp(-\alpha t))
      ,
   \end{equation}
   \MSRY{ADD!!}
\end{thm}

\begin{pf}
   See Theorem 2.4 in \cite{Tsukamoto:2021aa}.
\end{pf}


\MSRY{Any more comments, if needed}

\section{Problem Formulation}

Let us consider the following nonlinear system as follows:
\begin{equation} \label{eq:sys}
   \begin{aligned}   
      \ddtt{\mv{x}}
      =&
      \mv{f}(\mv{x},t)
      +
      \mm{B}(\mv{x},t)
      \mysat(\mv{u})   
      +
      \mv{d}(t)
      \\
      =&
      \mv{f}(\mv{x},t)
      +
      \mm{B}(\mv{x},t)
      \mv{u}
      +
      \mv{\phi(\mv{u})}
      +
      \mv{d}(t)
      ,
   \end{aligned}
\end{equation}
where $\mv{x}\in\R^n$ denotes the state variables, $\mv{u}\in\R^n$ denotes the control inputs, and $\mv{d}(t)\in\R^n$ denotes the external disturbances, and $\mv{f}:=\mv{f}(\mv{x},t)$ and $\mm{B}:=\mm{B}(\mv{x},t)$ denote system functions.
The input saturation function is denoted by $\mysat(\cdot)$ and is defined as
\begin{equation} \label{eq:sys:sat}
   \mysat(\mv{u}) =
   \begin{cases}
      \mv{u}, & \text{if } \norm{\mv{u}} \le \overline{u} \\
      \overline{u}\tfrac{\mv{u}}{\norm{\mv{u}}}, & \text{otherwise}
      ,
   \end{cases}
\end{equation}
where $\overline{u}\in\R_{>0}$ denotes the saturation threshold.
The difference between the controller output $\mv{u}$ and the saturated control input $\mysat(\mv{u})$ is denoted by $\mv{\phi}(\mv{u}) = \mysat(\mv{u}) - \mv{u}$.

In the practical perspective, we consider the following assumption:
\begin{assum}
   \cred{수도인버스로변경}The control input matrix $\mv{B}$ is sign-definite, so that the inverse of $\mm{B}$ exists.
\end{assum}

Motivated by the use of the partial contraction theory in \cite{Tsukamoto:2021af}, \MSRY{and more...}, we assume that the desired control input $\mv{u}_d\in\R^n$ which is designed to achieve the desired trajectory $\mv{x}_d\in\R^n$ under ideal conditions, \ie $\mv{d}=0$ and $\norm{\mv{u}_d}\le\overline{u}$, is given as follows:
\begin{equation} \label{eq:sys:desired}
   \ddtt\mv{x}_d   =
   \mv{f}_d
   +
   \mm{B}_d \mv{u}_d
   ,
\end{equation}
where $\mv{f}_d:=\mv{f}(\mv{x}_d, t)$ and $\mm{B}_d:=\mm{B}(\mv{x}_d,t)$.

Finally, the control objective is to design a control input $\mv{u}$ such that the system trajectories of the system \eqref{eq:pre:sys:perturbed} and the desired system \eqref{eq:sys:desired} converges, under the existences of the disturbances $\mv{d}$ and input saturation $\mysat(\cdot)$.

\section{Main Results}

\cred{
   \begin{itemize}
      \item 포화에 대한 영향을 강조하기 위한 내용 추가 (리마크)
      \item hard contraint를 쓰지 않고 LMI에 통합    
   \end{itemize}
}

Following the control design in \cite{Tsukamoto:2021af}, the state-dependent coefficient (SDC) formulation \cite{Dani:2015aa} is used to simplify the analysis to be understood as in LTV system theory.
The SDC $\mathbb{A}:=\mathbb{A}(\mv{x},\mv{x}_d)$ is defined as follows:
\begin{equation} \label{eq:SDC}
   \mathbb{A}
   :=
   \textstyle\int_0^1 
   \left[
      \pptfrac{\overline{\mv{f}}}{\mv{x}}     
   \right]
   \big(
      c \mv{x} + (1-c) \mv{x}_d
   \big)
   \,
   \der c
   ,
\end{equation}
where 
$
   \overline{\mv{f}}
   :=
   \overline{\mv{f}} (\mv{x},\mv{x}_d)
   =
   \mv{f}-\mv{f}_d
   + \mm{B} \mv{u}_d 
   - \mm{B}_d \mv{u}_d
$
and $c$ is a dummy variable for the integration.
The existence... blah blah can be found in \cite{Dani:2015aa}, \MSRY{and more...}.

Then, the error dynamics can be expressed as:
\begin{equation} \label{eq:error:dyn}
   \begin{aligned}
      \ddtt
      (
         \mv{x} - \mv{x}_d
      )
      =&
      \mv{f}-\mv{f}_d
      + \mm{B} \mv{u}
      - \mm{B}_d \mv{u}_d
      + \mm{B} \mv{\phi}(\mv{u}) + \mv{d}
      \\
      =&
      \underbrace
      {
         \mv{f}-\mv{f}_d
         +
         (
            \mm{B} 
            - \mm{B}_d 
         ) \mv{u}_d
      }_{= \mathbb{A}(\mv{x},\mv{x}_d) (\mv{x} - \mv{x}_d)}
      \\
      &
      + \mm{B} 
      (
         \mv{u} - \mv{u}_d
      )
      + \mm{B} \mv{\phi}(\mv{u}) + \mv{d}
      \\
      =&
      \mathbb{A} (\mv{x} - \mv{x}_d)
      + \mm{B} 
      (\mv{u} - \mv{u}_d)
      + \mm{B} \mv{\phi}(\mv{u}) + \mv{d}
      \\
      =&
      \left(
         \mathbb{A} - \mm{B} \mm{R}^{-1} \mm{B} ^\top \mm{M}
      \right)
      (\mv{x} - \mv{x}_d)
      + \mm{B} \mv{\phi}(\mv{u}) + \mv{d}
      ,
   \end{aligned}
\end{equation}
with the control input $\mv{u}$ defined as:
\begin{equation} \label{eq:control:input}
   \mv{u}
   =
   \mv{u}_d
   -
   \mm{R}^{-1} \mm{B}^\top \mm{M} (\mv{x} - \mv{x}_d)
\end{equation}
where $\mm{R}=\mm{R}^\top \succ 0$ is a symmetric positive definite matrix.

In the following, we formulate the optimization problem to obtain the contraction metric $\mm{M}$ which satisfies the contraction condition \eqref{eq:control:condition:contraction} and does not violate the input saturation constraint \eqref{eq:sys:sat}.


\subsection{Contraction Metric Calculation}

\subsubsection{Contraction Condition Satisfaction}

Invoking the partial contraction theory reviewed in Section \ref{sec:sub:contraction:partial}, we can define the virtual system as follows:
\begin{equation} \label{eq:sys:virtual1}
   \ddtt\mv{q}
   =
   \ddtt\mv{x}_d
   +
   \left(
      \mathbb{A} - \mm{B}\mm{R}^{-1} \mm{B}^\top \mm{M}
   \right)
   (\mv{q} - \mv{x}_d)
   +
   \mv{d}_{\mu}
   ,
\end{equation}
where $\mv{q}\in\R^n$ is the virtual state vector and $\mv{d}_{\mu}:=\mv{d}_{\mu}(\mv{d},\mu),\ \forall \mu\in[0,1]$, is defined as $\mv{0} = \mv{d}_{\mu}(\mv{d},0)$ and $\mv{d} = \mv{d}_{\mu}(\mv{d},1)$.
The differential dynamics is then defined as follows:
\begin{equation} \label{eq:sys:virtual2}
   \ddtt{\delta{\mv{q}}}
   =
   \left(
      \mathbb{A} - \mm{B}\mm{R}^{-1} \mm{B}^\top \mm{M}
   \right)
   \delta{\mv{q}}
   .
\end{equation}

\MSRY{CHECK FROM HERE}
Then according to Theorem \ref{thm:partial:contraction} and \ref{thm:contraction:perturbed}, system \eqref{eq:sys:virtual1} and the system \eqref{eq:sys:desired} will converge to the desired trajectory $\mv{x}_d$ with steady state error:
\begin{equation} \label{eq:ss:error}
   \mv{e}_{\rm{ss}}
   =
   \left(
      \sup({\mm{B}\mv{\phi(\mv{u})}})+\overline d
   \right)
   \tfrac{1}{\alpha}
   \sqrt{\tfrac{\overline m}{\underline{m}}}
   ,
\end{equation}
with the following conditions:
\begin{equation} \label{eq:control:condition:contraction}
   \ddtt\mm{M}
   +
   2\mysym(
      \mm{M}\mathbb{A}
   )
   +
   2\mm{M}\mm{B}\mm{R}^{-1} \mm{B}^\top \mm{M}
   \leq
   -2\alpha\mm{M}
   .
\end{equation}
By multiplying $\mm{W}:=\mm{M}^{-1}$ both sides of \eqref{eq:control:condition:contraction}, we have the following LMI:
\begin{equation} \label{eq:control:condition:contraction:LMI}
   -
   \ddtt\mm{W}
   +
   2\mysym(
      \mathbb{A}\mm{W}
   )
   -
   2\mm{B}\mm{R}^{-1} \mm{B}^\top 
   \leq
   -2\alpha\mm{W}
   .
\end{equation}

In \MSRY{[]}, Tsukamoto \etal proposed a method to formulate the LMI problem for the contraction metric as follows:
\begin{equation}
   \begin{aligned}
      &\min_{\mm{W},X,\mu} X + \lambda \mu
      \\
      &\text{s.t. \eqref{eq:control:condition:contraction}}
   \end{aligned}
\end{equation}

However, even with the proposed method, the input saturation can still occur by excessive contraction metric leading to potential instability.
Therefore, we need to incorporate input constraints into the LMI formulation.

\subsubsection{Input Saturation Constraint}

To address the input saturation function \eqref{eq:sys:sat}, the saturation constraint $\mv{u}^\top \mv{u} - \overline{u}^2 \le 0$ is defined as:
\begin{equation}
   \begin{bmatrix}
      \overline{u}^2 & \left(\mv{u}_d - \mm{R}^{-1}\mm{B}^\top\mm{M}(\mv{x}-\mv{x}_d)\right)^\top \\
      \star & \mm{I}
   \end{bmatrix}
   \succeq
   0
   .
\end{equation}
By multiplying both sides by $\left[\mydiag\right(\mm{I},\mm{W}(\mm{R}^{-1}\mm{B}^\top)^{-1}\left)\right]$, we have LMIs as follows:
\begin{equation} \label{eq:control:condition:saturation:1}
   \begin{bmatrix}
      \overline{u}^2 & \left(\mm{W}(\mm{R}^{-1}\mm{B}^\top)^{-1}\mv{u}_d - (\mv{x}-\mv{x}_d)\right)^\top \\
      \star & \mathbb{X}
   \end{bmatrix}
   \succeq
   0
   ,
\end{equation}
where $\mathbb{X}:=\mm{W}(\mm{R}^{-1}\mm{B}^\top)^{-1}{(\mm{R}^{-1}\mm{B}^\top)^{-1}}^\top\mm{W}$ which satisfies
\begin{equation} \label{eq:control:condition:saturation:2}
   \begin{bmatrix}
      \mathbb{X} & \mm{W}((\mm{R}^{-1}\mm{B}^\top)^{-1})^\top \\
      \star & \mm{I}
   \end{bmatrix}
   \succeq
   0
   .
\end{equation}

\hfill

Finally, the contraction metric can be obtained by solving the following LMI problem:
\begin{equation}
   \begin{aligned}
      &\min_{\mm{W},X,\mu} 
         X
         + 
         \rm{trace}(\mm{\mathbb{X}})
         \cred{
            + \lambda \mu
         }
      \\
      &\text{s.t. \eqref{eq:control:condition:contraction}, \eqref{eq:control:condition:saturation:1} and \eqref{eq:control:condition:saturation:2}}
   \end{aligned}
\end{equation}

\subsection{Numerical Methods For LMI}

\begin{itemize}
   \item Discretize $\ddtt\mv{M}$
   \item LMI solver
   \item Tuning methods
\end{itemize}


\section{Numerical Validation}

\subsection{Numerical Validation Set Up}

We employed the Lorenz system \MSRY{REF} to validate the proposed method which is represented as follows:
\begin{equation}
   \ddtt
   \underbrace{
      \begin{pmatrix}
         x\\y\\z
      \end{pmatrix}
   }_{=:\mv{\xi}}
   =
   \underbrace{
      \begin{pmatrix}
         \sigma(y-x)
         \\
         x(\rho-z)-y
         \\
         xy-\beta z
      \end{pmatrix}
   }_{=:\mv{f}(\mv{\xi})}
   +
   \underbrace{
      \begin{pmatrix}
         u_x\\u_y\\u_z
      \end{pmatrix}
   }_{=:\mv{u}}
   +
   \underbrace{
      \begin{pmatrix}
         d_x\\d_y\\d_z
      \end{pmatrix}
   }_{=:\mv{d}}
   ,
\end{equation}
where $\sigma=10$, $\rho=28$, and $\beta=\tfrac{8}{3}$ are system parameters.



\subsection{Numerical Validation Results}

\begin{itemize}
   \item \colorLine{black}{solid} Desired 
   \item \colorLine{blue}{dashed} CMP1 
   \item \colorLine{red}{dashed} CMP2 
   \item \colorLine{cyan}{dashed} CMP3 
\end{itemize}


\begin{figure*}[t]
   \centering
   \subfloat[State trajectories.]{
      \includegraphics[width=0.49\linewidth]
      {
      src/script_simulation/figures/compare/Fig1.eps
      }%
      \label{fig:result:state}}
   \hfill
   \subfloat[Control inputs.]{
      \includegraphics[width=0.49\linewidth]
      {
      src/script_simulation/figures/compare/Fig3.eps
      }%
      \label{fig:result:control}}
   \vfill
   \subfloat[Norms of tracking errors.]{
      \includegraphics[width=0.49\linewidth]
      {
      src/script_simulation/figures/compare/Fig6.eps
      }%
      \label{fig:result:error:norm}}
   \hfill
   \subfloat[Norms of control inputs.]{
      \includegraphics[width=0.49\linewidth]
      {
      src/script_simulation/figures/compare/Fig5.eps
      }%
      \label{fig:result:control:norm}}
   \caption{
   Numerical validation results of (C$_1$) \protect\colorLine{blue}{solid}, (C$_2$) \protect\colorLine{cyan}{solid}, (C$_3$) \protect\colorLine{gray}{solid}, (C$_4$) \protect\colorLine{magenta}{solid}, and desired trajectory $x_d$ \protect\colorLine{red}{dashed}.
}
\label{fig:result}
\end{figure*}

\begin{figure}[t]
   \centering
   \subfloat[State trajectories (3D-view).]{
      \includegraphics[width=0.48\linewidth]
      {
      src/script_simulation/figures/compare/Fig2.eps
      }%
      \label{fig:result:state:3d}}
   \hfill
   \subfloat[Control input loci (3D-view).]{
      \includegraphics[width=0.48\linewidth]
      {
      % src/script_simulation/figures/compare/Fig4.eps
      src/script_simulation/figures/compare/Fig4.png
      }%
      \label{fig:result:control:3d}}
   \caption{
   Numerical validation results of (C$_1$) \protect\colorLine{blue}{solid}, (C$_2$) \protect\colorLine{cyan}{solid}, (C$_3$) \protect\colorLine{gray}{solid}, (C$_4$) \protect\colorLine{magenta}{solid}, and desired trajectory $x_d$ \protect\colorLine{red}{dashed}.
}
\label{fig:result:3D}
\end{figure}

\section{Conclusion}

In this study, ...

\begin{comment}

\section{Introduction}
This document is a template for \LaTeXe. If you are reading a paper or
PDF version of this document, please download the electronic file
\texttt{ifacconf.tex}. You will also need the class file
\texttt{ifacconf.cls}. Both files are available on the IFAC web site.

Please stick to the format defined by the \texttt{ifacconf} class, and
do not change the margins or the general layout of the paper. It
is especially important that you do not put any running header/footer
or page number in the submitted paper.\footnote{
This is the default for the provided class file.}
Use \emph{italics} for emphasis; do not underline.

Page limits may vary from conference to conference. Please observe the 
page limits of the event for which your paper is intended.

% $\vec$

\section{Procedure for Paper Submission}

Next we see a few subsections.

\subsection{Review Stage}

For submission guidelines, follow instructions on paper submission
system as well as the event website.

Note that conferences impose strict page limits, so it will be better
for you to prepare your initial submission in the camera ready layout
so that you will have a good estimate for the paper
length. Additionally, the effort required for final submission will be
minimal.

\subsection{Equations}

Some words might be appropriate describing equation~(\ref{eq:sample}), if 
we had but time and space enough. 

\begin{equation} \label{eq:sample}
{{\partial F}\over {\partial t}} = D{{\partial^2 F}\over {\partial x^2}}.
\end{equation}

See \cite{Abl:56}, \cite{AbTaRu:54}, \cite{Keo:58} and \cite{Pow:85}.

\subsubsection{Example.} This equation goes far beyond the
celebrated theorem ascribed to the great Pythagoras by his followers.

\begin{thm}   % use the thm environment for theorems
The square of the length of the hypotenuse of a right triangle equals
the sum of the squares of the lengths of the other two sides.
\end{thm}

\begin{pf}    % and the pf environment for proofs
The square of the length of the hypotenuse of a right triangle equals the sum of the squares 
of the lengths of the other two sides.
\end{pf}

%% There are a number of predefined theorem-like environments in
%% ifacconf.cls:
%%
%% \begin{thm} ... \end{thm}            % Theorem
%% \begin{lem} ... \end{lem}            % Lemma
%% \begin{claim} ... \end{claim}        % Claim
%% \begin{conj} ... \end{conj}          % Conjecture
%% \begin{cor} ... \end{cor}            % Corollary
%% \begin{fact} ... \end{fact}          % Fact
%% \begin{hypo} ... \end{hypo}          % Hypothesis
%% \begin{prop} ... \end{prop}          % Proposition
%% \begin{crit} ... \end{crit}          % Criterion

Of course LaTeX manages equations through built-in macros. You may
wish to use the \texttt{amstex} package for enhanced math
capabilities.

\subsection{Figures}

To insert figures, use the \texttt{graphicx} package. Although other
graphics packages can also be used, \texttt{graphicx} is simpler to
use. See  Fig.~\ref{fig:bifurcation} for an example.

\begin{figure}
\begin{center}
\includegraphics[width=8.4cm]{bifurcation}    % The printed column width is 8.4 cm.
\caption{Bifurcation: Plot of local maxima of $x$ with damping $a$ decreasing} 
\label{fig:bifurcation}
\end{center}
\end{figure}

Figures must be centered, and have a caption at the bottom. 

\subsection{Tables}
Tables must be centered and have a caption above them, numbered with
Arabic numerals. See table~\ref{tb:margins} for an example.

\begin{table}[hb]
\begin{center}
\caption{Margin settings}\label{tb:margins}
\begin{tabular}{cccc}
Page & Top & Bottom & Left/Right \\\hline
First & 3.5 & 2.5 & 1.5 \\
Rest & 2.5 & 2.5 & 1.5 \\ \hline
\end{tabular}
\end{center}
\end{table}

\subsection{Final Stage}

Authors are expected to mind the margins diligently.  Papers need to
be stamped with event data and paginated for inclusion in the
proceedings. If your manuscript bleeds into margins, you will be
required to resubmit and delay the proceedings preparation in the
process.

\subsubsection{Page margins.} See table~\ref{tb:margins} for the
page margins specification. All dimensions are in \emph{centimeters}.


\subsection{PDF Creation}

All fonts must be embedded/subsetted in the PDF file. Use one of the
following tools to produce a good quality PDF file:

\subsubsection{PDFLaTeX} is a special version of LaTeX by Han The
Thanh which produces PDF output directly using Type-1 fonts instead of
the standard \texttt{dvi} file. It accepts figures in JPEG, PNG, and PDF
formats, but not PostScript. Encapsulated PostScript figures can be
converted to PDF with the \texttt{epstopdf} tool or with Adobe Acrobat
Distiller.

\subsubsection{Generating PDF from PostScript} is the classical way of
producing PDF files from LaTeX. The steps are:

\begin{enumerate}
  \item Produce a \texttt{dvi} file by running \texttt{latex} twice.
  \item Produce a PostScript (\texttt{ps}) file with \texttt{dvips}.
  \item Produce a PDF file with \texttt{ps2pdf} or Adobe Acrobat
  Distiller.
\end{enumerate}

\subsection{Copyright Form}

IFAC will put in place an electronic copyright transfer system in due
course. Please \emph{do not} send copyright forms by mail or fax. More
information on this will be made available on IFAC website.


\section{Units}

Use SI as primary units. Other units may be used as secondary units
(in parentheses). This applies to papers in data storage. For example,
write ``$15\,\mathrm{Gb}/\mathrm{cm}^2$ ($100\,\mathrm{Gb}/\mathrm{in}^2$)''. 
An exception is when
English units are used as identifiers in trade, such as ``3.5 in
disk drive''. Avoid combining SI and other units, such as current in
amperes and magnetic field in oersteds. This often leads to confusion
because equations do not balance dimensionally. If you must use mixed
units, clearly state the units for each quantity in an equation.  The
SI unit for magnetic field strength $\mathbf{H}$ is $\mathrm{A}/\mathrm{m}$. However, if you wish to
use units of $\mathrm{T}$, either refer to magnetic flux density $\mathbf{B}$ or
magnetic field strength symbolized as $\mu_0\,\mathbf{H}$. Use the center dot to
separate compound units, e.g., ``$\mathrm{A} \cdot \mathrm{m}^2$''.

\section{Helpful Hints}

\subsection{Figures and Tables}

Figure axis labels are often a source of confusion. Use words rather
than symbols. As an example, write the quantity ``Magnetization'', or
``Magnetization M'', not just ``M''. Put units in parentheses. Do not
label axes only with units.  For example, write ``Magnetization
($\mathrm{A}/\mathrm{m}$)'' or ``Magnetization ($\mathrm{A} \mathrm{m}^{-1}$)'', not just
 ``$\mathrm{A}/\mathrm{m}$''. Do not
label axes with a ratio of quantities and units. For example, write
``Temperature ($\mathrm{K}$)'', not ``$\mbox{Temperature}/\mathrm{K}$''.

Multipliers can be especially confusing. Write ``Magnetization
($\mathrm{kA}/\mathrm{m}$)'' or ``Magnetization ($10^3 \mathrm{A}/\mathrm{m}$)''. Do not write
``Magnetization $(\mathrm{A}/\mathrm{m}) \times 1000$'' because the reader would not know
whether the axis label means $16000\,\mathrm{A}/\mathrm{m}$ or $0.016\,\mathrm{A}/\mathrm{m}$.

\subsection{References}

Use Harvard style references (see at the end of this document). With
\LaTeX, you can process an external bibliography database 
using \texttt{bibtex},\footnote{In this case you will also need the \texttt{ifacconf.bst}
file, which is part of the \texttt{ifaconf} package.}
or insert it directly into the reference section. Footnotes should be avoided as
far as possible.  Please note that the references at the end of this
document are in the preferred referencing style. Papers that have not
been published should be cited as ``unpublished''.  Capitalize only the
first word in a paper title, except for proper nouns and element
symbols.

\subsection{Abbreviations and Acronyms}

Define abbreviations and acronyms the first time they are used in the
text, even after they have already been defined in the
abstract. Abbreviations such as IFAC, SI, ac, and dc do not have to be
defined. Abbreviations that incorporate periods should not have
spaces: write ``C.N.R.S.'', not ``C. N. R. S.'' Do not use abbreviations
in the title unless they are unavoidable (for example, ``IFAC'' in the
title of this article).

\subsection{Equations}

Number equations consecutively with equation numbers in parentheses
flush with the right margin, as in (\ref{eq:sample}).  To make your equations more
compact, you may use the solidus ($/$), the $\exp$ function, or
appropriate exponents. Use parentheses to avoid ambiguities in
denominators. Punctuate equations when they are part of a sentence, as
in

\begin{equation} \label{eq:sample2}
\begin{array}{ll}
\int_0^{r_2} & F (r, \varphi ) dr d\varphi = [\sigma r_2 / (2 \mu_0 )] \\
& \cdot \int_0^{\inf} exp(-\lambda |z_j - z_i |) \lambda^{-1} J_1 (\lambda  r_2 ) J_0 (\lambda r_i ) d\lambda 
\end{array}
\end{equation}

Be sure that the symbols in your equation have been defined before the
equation appears or immediately following. Italicize symbols ($T$
might refer to temperature, but T is the unit tesla). Refer to
``(\ref{eq:sample})'', not ``Eq. (\ref{eq:sample})'' or ``equation
(\ref{eq:sample})'', except at the beginning of a sentence: ``Equation
(\ref{eq:sample}) is \ldots''.

\subsection{Other Recommendations}

Use one space after periods and colons. Hyphenate complex modifiers:
``zero-field-cooled magnetization''. Avoid dangling participles, such
as, ``Using (1), the potential was calculated'' (it is not clear who or
what used (1)). Write instead: ``The potential was calculated by using
(1)'', or ``Using (1), we calculated the potential''.

A parenthetical statement at the end of a sentence is punctuated
outside of the closing parenthesis (like this). (A parenthetical
sentence is punctuated within the parentheses.) Avoid contractions;
for example, write ``do not'' instead of ``don' t''. The serial comma
is preferred: ``A, B, and C'' instead of ``A, B and C''.


\section{Conclusion}

A conclusion section is not required. Although a conclusion may review
the main points of the paper, do not replicate the abstract as the
conclusion. A conclusion might elaborate on the importance of the work
or suggest applications and extensions.

\end{comment}

\begin{ack}
Place acknowledgments here.
\end{ack}

\bibliography{template/refs, localRefs}             % bib file to produce the bibliography
                                    % with bibtex (preferred)

% \appendix
% \section{SDC Derivation} \label{sec:appendix:SDC}

% In this appendix, the SDC $\mathbb{A}$ will be derived in detail.
% First, we follow the definition of the SDC in \eqref{eq:SDC} as follows:
% \begin{equation}
%    \begin{aligned}
%       \mathbb{A}
%       =&
%       \textstyle\int_0^1 
%       \left[
%          \pptfrac{}{\mv{x}}
%          \left(
%             \mm{A} \mv{x}
%             +
%             \mm{B} \mv{u}_d
%          \right)
%       \right]
%       \big(
%          c \mv{x} + (1-c) \mv{x}_d
%       \big)
%       \,
%       \der c
%       \\
%       =&
%       \textstyle\int_0^1
%       \left[
%          \left(
%             \mm{A}
%             +
%             \begin{bmatrix}
%                0 & 0 
%                \\
%                b_1 & 0
%             \end{bmatrix}
%          \right)
%       \right]
%       \big(
%          c \mv{x} + (1-c) \mv{x}_d      
%       \big)
%       \,
%       \der c
%       ,
%    \end{aligned}
% \end{equation}
% where 
% $
%    b_1
%    :=
%    b_1(\mv{x},\mv{x}_d)
%    = 
%    - \tfrac{3P^2\Phi}{2J}
%    \left(
%       \cos(P x_1) {u_1}_d
%       +
%       \sin(P x_1) {u_2}_d
%    \right)
% $.
% Then, we can rewrite the SDC as follows:
% \begin{equation} \label{eq:SDC:explicit:2}
%    \mathbb{A}
%    =
%    \begin{bmatrix}
%       0 & 1 \\
%       b_2(\mv{x},\mv{x}_d) & -\tfrac{f_v}{J}
%    \end{bmatrix}
%    ,
% \end{equation}
% where 
% \begin{equation*}
%    \begin{aligned}
%       b_2(\mv{x},\mv{x}_d)
%       :=&
%       \tfrac{3P\Phi}{2J(x_1-{x_1}_d)}
%       \big[
%          \left\{
%             \sin(P x_1)
%             -
%             \sin(P{ x_1}_d)
%          \right\}
%          {u_1}_d
%       \\
%       &
%          +
%          \left\{
%             \cos(P x_1)
%             -
%             \cos(P {x_1}_d)
%          \right\}
%          {u_2}_d
%       \big]
%       .
%    \end{aligned}
% \end{equation*}

\end{document}
